\documentclass{ctexart}

% 数学相关宏包
\usepackage{amsmath,mathtools,amsthm,amssymb}
\usepackage{bbm}
\usepackage{mathrsfs}
\usepackage{centernot}

% 图形和绘图
\usepackage{graphicx,float}
\usepackage{tikz,tikz-cd}
\usepackage{pgfplots}
\usepackage[framemethod=TikZ]{mdframed}
\usetikzlibrary{decorations.pathmorphing,positioning,arrows,automata,calc,shapes,patterns}
\pgfplotsset{compat=1.18}
\usepgfplotslibrary{statistics}

% 表格和列表
\usepackage{booktabs,multirow,colortbl}
\usepackage{enumitem}
\setlist[itemize,enumerate,description]{noitemsep}
\setlistdepth{9}

% 算法
\usepackage{algorithm,algorithmic}

% 字体和颜色
\usepackage{xcolor,listings}

% 页面和链接
\usepackage{fancyhdr,caption}
\usepackage{hyperref}
\usepackage{cleveref}

% 工具宏包
\usepackage{etoolbox,letltxmacro,docmute}
\usepackage{gensymb}

% 自定义设置
\renewcommand{\arraystretch}{1.5}
\let\originalmathbb\mathbb
\renewcommand{\mathbb}[1]{%
    \ifstrequal{#1}{1}{\mathbbm{#1}}{\originalmathbb{#1}}%
}

% 颜色定义
\definecolor{lightblue}{RGB}{173,216,230}
\definecolor{lightgreen}{RGB}{144,238,144}
\definecolor{lightyellow}{RGB}{255,255,224}
\definecolor{lightgray}{RGB}{211,211,211}

% 超链接设置
\hypersetup{
    colorlinks=true,
    linkcolor=red,
    filecolor=magenta,
    urlcolor=cyan,
    pdftitle={数学笔记},
    pdfauthor={乐绎华},
    pdfsubject={数学笔记},
    pdfkeywords={数学},
    bookmarksnumbered=true,
    bookmarksopen=true,
    bookmarksopenlevel=6,
    pdfstartview=FitH,
    pdfpagemode=UseOutlines,
    pdfborder={0 0 0}
}

% 定理环境
\theoremstyle{plain}
\newtheorem{theorem}{定理}[section]
\newtheorem{lemma}[theorem]{引理}
\newtheorem{corollary}[theorem]{推论}
\newtheorem{proposition}[theorem]{命题}

\theoremstyle{definition}
\newtheorem{definition}[theorem]{定义}
\newtheorem{example}[theorem]{例}
\newtheorem{exercise}[theorem]{练习}
\newtheorem{problem}[theorem]{题目}

\theoremstyle{remark}
\newtheorem{remark}[theorem]{备注}
\newtheorem{note}[theorem]{注}
\newtheorem{observation}[theorem]{观察}

% 解答环境
\newtheorem*{solution}{解}
\newtheorem*{proof*}{证明}

% 图片路径
\graphicspath{{F:/iCloudDrive/iCloud~md~obsidian/2025-summer/figures/}}

\author{乐绎华}
\title{数学笔记}
\date{\today}

\begin{document}

\maketitle
\tableofcontents

\href{https://moroianu.perso.math.cnrs.fr/tex/kg.pdf}{kg.dvi}

Examples:
\begin{figure}[H]
\centering
\includegraphics[width=\textwidth]{1-kahler-geometry-2025082312.png}
% \caption{}
\label{}
\end{figure}

Applications:
\begin{figure}[H]
\centering
\includegraphics[width=\textwidth]{1-kahler-geometry-2025082313.png}
% \caption{}
\label{}
\end{figure}

\begin{figure}[H]
\centering
\includegraphics[width=\textwidth]{kahler-geometry-2025082312.png}
% \caption{}
\label{}
\end{figure}

The motivation of the definition of $J$ is from $j$, the multiplication by $i$ on $\mathbb{C}$ via the identification of $\mathbb{R}^{2}$ with $\mathbb{C}$ given by $z=x+iy\mapsto(x,y)$. The endomorphism $j$ can be expressed in the canonical base as
\[
j=\begin{pmatrix}
0 & -1 \\
1 & 0
\end{pmatrix}
\]
The differential of $F=f+ig:U\subset \mathbb{C}\to \mathbb{C}$ (viewed as a real function $F:U\subset \mathbb{R}^{2}\to \mathbb{R}^{2}$) is the linear map
\[
F_{*}=\begin{pmatrix}
\frac{ \partial f }{ \partial x }  & \frac{ \partial f }{ \partial y }  \\
\frac{ \partial g }{ \partial x }  & \frac{ \partial g }{ \partial y } 
\end{pmatrix}
\]
Then it is easy to check that the Cauchy-Riemann relations are equivalent to the commutation relation $jF_{*}=F_{*}j$.

Then for $F:U\subset \mathbb{C}^{n}\to \mathbb{C}^{m}$, it is holomorphic iff the diffrerential $F_{*}$ of $F$ as real map $F:\mathbb{R}^{2n}\to \mathbb{R}^{2m}$ satisfies $j_mF_{*}=F_{*}j_n$, where
\[
j_m=\begin{pmatrix}
0 & -I _m \\
I _m & 0
\end{pmatrix}
\]
\begin{figure}[H]
\centering
\includegraphics[width=\textwidth]{kahler-geometry-2025082313.png}
% \caption{}
\label{}
\end{figure}
\begin{figure}[H]
\centering
\includegraphics[width=\textwidth]{kahler-geometry-2025082314.png}
% \caption{}
\label{}
\end{figure}

\begin{figure}[H]
\centering
\includegraphics[width=\textwidth]{1-kahler-geometry-2025082314.png}
% \caption{}
\label{}
\end{figure}

Almost complex structure $J$ :
\begin{figure}[H]
\centering
\includegraphics[width=\textwidth]{2-kahler-geometry-2025082314.png}
% \caption{}
\label{}
\end{figure}
$J_{U}$ does not depend on $U$ and their collection is a well-defined tensor $J$ on $M$. This tensor clearly satisfies $J^{2}=-id$.

\begin{figure}[H]
\centering
\includegraphics[width=\textwidth]{3-kahler-geometry-2025082314.png}
% \caption{}
\label{}
\end{figure}
A $(1,1)$ -tensor $J$ is a type of tensor that takes in one vector and one covector (a linear function) and outputs a scalar.

A complex manifold is thus in a canonical way an almost complex manifold.

\begin{figure}[H]
\centering
\includegraphics[width=\textwidth]{4-kahler-geometry-2025082314.png}
% \caption{}
\label{}
\end{figure}
\begin{figure}[H]
\centering
\includegraphics[width=\textwidth]{5-kahler-geometry-2025082314.png}
% \caption{}
\label{}
\end{figure}

\textbf{Integrability} of the distribution $T^{0,1}M$ means that its sections are closed under the Lie bracket, which is a condition that ensires the existence of a compatible complex coordinate system. In other words:
\[
[X,Y]\in \Gamma(D)\qquad \forall X,Y\in\Gamma(D)
\]
where the distribution $D$ is a subbundle of $T_{\mathbb{C}}M$ and $\Gamma(D)$ is a vector field in $D$.


Suppose that $J$ comes from a holomorphic structure on $M$. Consider a local chart $(U,\phi_{U})$ and let $z_{\alpha}=x_{\alpha}+iy_{\alpha}$ be the $\alpha$ -th component of $\phi_{U}$. If $\{ e_1,\dots,e_{2m} \}$ denotes the standard basis of $\mathbb{R}^{2m}$, we have by definition:
\[
\frac{ \partial   }{ \partial x_{\alpha} } =(\phi_{U})^{-1}_{*}(e_{\alpha}),\qquad \frac{ \partial   }{ \partial y_{\alpha} } =(\phi_{U})^{-1}_{*}(e_{m+\alpha})
\]
Moreover, $j_m(e_{\alpha})=e_{m+\alpha}$, so we obtain directly from the defintion
\[
J\left( \frac{ \partial   }{ \partial x_{\alpha} }  \right)=\frac{ \partial   }{ \partial y_{\alpha} }
\]
We now make the following notations
\[
\frac{\partial}{\partial z_\alpha}:=\frac{1}{2}\left(\frac{\partial}{\partial x_\alpha}-i \frac{\partial}{\partial y_\alpha}\right), \quad \frac{\partial}{\partial \bar{z}_\alpha}:=\frac{1}{2}\left(\frac{\partial}{\partial x_\alpha}+i \frac{\partial}{\partial y_\alpha}\right) .
\]
Suppose that $u_\alpha, v_\alpha$ is another local system of coordinates, satisfying
\[
\frac{\partial}{\partial v_\alpha}=J \frac{\partial}{\partial u_\alpha}
\]
We then have
\[
\frac{\partial u_\beta}{\partial x_\alpha}=\frac{\partial v_\beta}{\partial y_\alpha} \quad \text { and } \quad \frac{\partial u_\beta}{\partial y_\alpha}=-\frac{\partial v_\beta}{\partial x_\alpha},
\]
thus the transition functions are holomorphic.

\begin{exercise}
\begin{figure}[H]
\centering
\includegraphics[width=\textwidth]{kahler-geometry-2025082711.png}
% \caption{}
\label{}
\end{figure}\label{7c5a9f}
\end{exercise}

\begin{proof}
\[
\dots=\begin{pmatrix}
A+iB & B \\
O & A-iB
\end{pmatrix}
\]
Then
\[
\det \begin{pmatrix}
A & B \\
-B & A
\end{pmatrix}=\det (A+iB)\cdot \det(A-iB)=\det(A+iB)\cdot  \overline{\det(A+iB)}=\lvert \det(A+iB) \rvert ^{2}>0
\]
\end{proof}

\begin{exercise}
\begin{figure}[H]
\centering
\includegraphics[width=\textwidth]{1-kahler-geometry-2025082711.png}
% \caption{}
\label{}
\end{figure}
\end{exercise}
A smooth manifold is \textbf{orientable} if you can find an \textbf{oriented atlas}.

Given a collection of coordinate charts $(U_k,\phi _k)$ that covers the manifold, we need to show that the Jacobian determinant of all transition functions $\phi _j\circ \phi _k^{-1}$ is everywhere positive.

Under the basis $(v_1,\dots,v_n,Jv_1,\dots,Jv_n)$, any transition functions have the real matrix form
\[
\begin{pmatrix}
A & B \\
-B & A
\end{pmatrix}
\]
which has strictly positive determinant by \cref{7c5a9f} .

\section{Holomorphic forms and vector fields}

\begin{figure}[H]
\centering
\includegraphics[width=\textwidth]{6-kahler-geometry-2025082314.png}
% \caption{}
\label{}
\end{figure}

\begin{figure}[H]
\centering
\includegraphics[width=\textwidth]{kahler-geometry-2025082315.png}
% \caption{}
\label{}
\end{figure}
\begin{figure}[H]
\centering
\includegraphics[width=\textwidth]{1-kahler-geometry-2025082315.png}
% \caption{}
\label{}
\end{figure}

\begin{figure}[H]
\centering
\includegraphics[width=\textwidth]{kahler-geometry-2025082713.png}
% \caption{}
\label{}
\end{figure}
$[\cdot,\cdot]$ is the Lie Bracket, which is defined by
\[
[X,Y](f)=X(Y(f))-Y(X(f))
\]
\begin{proposition}
\begin{figure}[H]
\centering
\includegraphics[width=\textwidth]{1-kahler-geometry-2025082713.png}
% \caption{}
\label{}
\end{figure}
\begin{figure}[H]
\centering
\includegraphics[width=\textwidth]{2-kahler-geometry-2025082713.png}
% \caption{}
\label{}
\end{figure}
\end{proposition}
\begin{proof}
Read.
\end{proof}

\begin{figure}[H]
\centering
\includegraphics[width=\textwidth]{3-kahler-geometry-2025082713.png}
% \caption{}
\label{}
\end{figure}
\begin{figure}[H]
\centering
\includegraphics[width=\textwidth]{kahler-geometry-2025082714.png}
% \caption{}
\label{}
\end{figure}
\begin{figure}[H]
\centering
\includegraphics[width=\textwidth]{1-kahler-geometry-2025082714.png}
% \caption{}
\label{}
\end{figure}
\begin{figure}[H]
\centering
\includegraphics[width=\textwidth]{kahler-geometry-2025082715.png}
% \caption{}
\label{}
\end{figure}

\section{Hermitian bundles}

Let $M$ be a differentiable manifold (not necessarily complex) and let $E\to M$ be a \underline{complex vector bundle} over $M$, which means $E$ is a space with a projection map to $M$, and the fibers over each point in $M$ are complex vector space.

\begin{definition}[connection]
A ($\mathbb{C}$ -linear) connection on a $E$ is a $\mathbb{C}$ -linear differential operator $\nabla:C^{\infty}(E)\to C^{\infty}(\Lambda^{1}(E))$ satisfying the Leibniz rule
\[
\nabla(f\sigma)=df\otimes \sigma+f\nabla\sigma,\qquad \forall f\in C^{\infty}(M)
\]
\end{definition}
$C^{\infty}(\Lambda^{1}(E))$ is the space of $E$ -valued 1-forms, where $\Lambda^{1}(E)$ denotes the bundle of $E$ -valued 1-forms over $M$.

One can extend any connection to the bundles of $E$ -valued $p$ -forms on $M$ by
\[
\nabla(\omega \otimes \sigma)=d\omega \otimes \sigma+(-1)^{p}\omega \wedge \nabla\sigma
\]
where the wedge product has to be understood as
\[
\omega \wedge \nabla\sigma\coloneqq \sum_{i=1}^{n} \omega \wedge e^{ * }_i\otimes \nabla_{e_i}\sigma
\]
for any local basis $\{ e_i \}$ of $TM$ with dual basis $\{ e^{ * }_i \}$. $\nabla_{e_i}\sigma$ is the covariant derivative along the basis vector $e_i$.

\begin{definition}[curvature]
The \textbf{curvature} of $\nabla$ is the $\text{End}(E)$ -valued 2-form $R^{\nabla}$ defined by
\[
R^{\nabla}(\sigma)\coloneqq \nabla(\nabla\sigma)\qquad \forall \sigma\in C^{\infty}(E)
\]
\end{definition}
To check that this is indeed tensorial (i.e. depends on $\sigma$ and not on $f$), we can write
\[
\nabla^{2}(f\sigma)=\nabla(df\otimes \sigma+f\nabla\sigma)=\underbrace{ d^{2}f }_{ =0 }\otimes \sigma-df\wedge \nabla\sigma+df\wedge \nabla\sigma+f\nabla^{2}\sigma=f\nabla^{2}\sigma
\]
\begin{figure}[H]
\centering
\includegraphics[width=\textwidth]{kahler-geometry-2025082819.png}
% \caption{}
\label{}
\end{figure}
\begin{figure}[H]
\centering
\includegraphics[width=\textwidth]{1-kahler-geometry-2025082819.png}
% \caption{}
\label{}
\end{figure}
\begin{figure}[H]
\centering
\includegraphics[width=\textwidth]{2-kahler-geometry-2025082819.png}
% \caption{}
\label{}
\end{figure}


\end{document}